% -*- coding: utf-8 -*-
% !TEX program = xelatex
\documentclass[plain,noanswer,medmath]{../../template/jnuexam}
\usepackage{../../template/kysx}

\begin{document}


\examtitle{2016}{3}


\loadproblems{1}{2016P1}%载入试卷一
\loadproblems{2}{2016P2}%载入试卷二


\makepart{选择题}{1～8小题，每小题4分，共32分}


\useproblem{2}{1}{4}%【同试卷二第一(4)题】


\useproblem{2}{1}{6}%【同试卷二第一(6)题】


\begin{problem}
设$J_i = \iint_{D_i}\sqrt[3]{x-y}\dx\dy$（$i=1,2,3$），其中
\begin{align*}
  D_1 &  = \left\{(x,y) \middle| 0 \le x \le 1,0 \le y \le 1 \right\}, \\
  D_2 &  = \left\{(x,y) \middle| 0 \le x \le 1,0 \le y \le \sqrt{x} \right\}, \\
  D_3 &  = \left\{(x,y) \middle| 0 \le x \le 1,x^2 \le y \le 1 \right\}.
\end{align*}
则\pickout{}
\begin{abcd}
\item $J_1 < J_2 < J_3$．
\item $J_3 < J_1 < J_2$．
\item $J_2 < J_3 < J_1$．
\item $J_2 < J_1 < J_3$．
\end{abcd}
\end{problem}


\begin{problem}
级数为$\sum_{n = 1}^{\infty}\left(\frac{1}{\sqrt{n}} - \frac{1}{\sqrt{n + 1}} \right)\sin(n + k)$%
（$k$为常数），则该级数\pickout{}
\begin{abcd}
\item 绝对收敛．
\item 条件收敛．
\item 发散．
\item 收敛性与$k$有关．
\end{abcd}
\end{problem}


\useproblem{1}{1}{5}%【同试卷一第一(5)题】


\useproblem{2}{1}{8}%【同试卷二第一(8)题】


\begin{problem}
设$A,B$为随机事件，$0 < P(A) < 1$，$0 < P(B) < 1$．若$P(A\left| B \right.) = 1$，则下面正确的是\pickout{}
\begin{abcd}
\item $P(\bar B\left| \bar A \right.) = 1$．
\item $P(A\left| \bar B \right.) = 0$．
\item $P(A + B) = 1$．
\item $P(B\left| A \right.) = 1$．
\end{abcd}
\end{problem}


\begin{problem}
设随机变量$X,Y$独立，且$X\sim N(1,2)$，$Y\sim (1,4)$，则$D(XY)$为\pickout{}
\begin{abcd}
\item $6$．
\item $8$．
\item $14$．
\item $15$．
\end{abcd}
\end{problem}


\makepart{填空题}{9～14小题，每小题4分，共24分}


\begin{problem}
已知函数$f(x)$满足$\lim_{x \to 0} \frac{\sqrt{1 + f(x)\sin 2x} - 1}{\e^{3x} - 1} = 2$，
则$\lim_{x \to 0} f(x) =$ \fillin{}．
\end{problem}


\useproblem{2}{2}{10}%【同试卷二第二(10)题】


\useproblem{1}{2}{11}%【同试卷一第二(11)题】


\begin{problem}
设$D = \left\{(x,y)\;\middle|\;|x| \le y \le 1, - 1 \le x \le 1 \right\}$，
则$\iint_D x^2\e^{- y^2} \dx\dy =$ \fillin{}．
\end{problem}


\useproblem{1}{2}{13}%【同试卷一第二(13)题】


\begin{problem}
设袋中有红、白、黑球各$1$个，从中有放回的取球，每次取$1$个，直到三种颜色的球都取到为止，
则取球次数恰为$4$的概率为 \fillin{}．
\end{problem}


\makepart{解答题}{15～23小题，共94分}


\useproblem[points=10]{2}{3}{15}%【同试卷二第三(15)题】


\begin{problem}[points=10]
设某商品的最大需求量为$1200$件，该商品的需求函数$Q = Q(p)$，
需求弹性$\eta = \frac{p}{120 - p}$（$\eta > 0$），$p$为单价（万元）．
\begin{enumerate}
\item 求需求函数的表达式；
\item 求$p = 100$万元时的边际收益，并说明其经济意义．
\end{enumerate}
\end{problem}


\useproblem[points=10]{2}{3}{16}%【同试卷二第三(16)题】


\begin{problem}[points=10]
设函数$f(x)$连续，且满足$\int_0^x f(x - t) \dt = \int_0^x (x - t)f(t)\dt + \e^{- x} - 1$，求$f(x)$．
\end{problem}


\begin{problem}[points=10]
求幂级数$\sum_{n = 0}^{\infty}\frac{x^{2n + 2}}{(n + 1)(2n + 1)}$的收敛域及和函数．
\end{problem}


\useproblem[points=11]{2}{3}{22}%【同试卷二第三(22)题】


\useproblem[points=11]{1}{3}{21}%【同试卷一第三(21)题】


\useproblem[points=11]{1}{3}{22}%【同试卷一第三(22)题】


\useproblem[points=11]{1}{3}{23}%【同试卷一第三(23)题】


\end{document}
